\chapter{Definições, características e aplicações}
\label{cap:definicoes}

\begin{objetivos}
Ao final deste capítulo você deve ser capaz de:
\begin{itemize}
  \item definir CLP a partir da definição normativa e de uma definição operacional;
  \item reconhecer as características que qualificam um equipamento como controlador industrial;
  \item identificar os blocos que compõem um CLP e o que cada um faz;
  \item diferenciar controlador compacto, modular e com I/O distribuído;
  \item distinguir o CLP de um computador de uso geral.
\end{itemize}
\end{objetivos}

\section{Definição}

A definição clássica de controlador programável, dada pela \emph{National Electrical Manufacturers
Association} (NEMA), é:

\begin{quote}
\emph{``Aparelho eletrônico digital que utiliza uma memória programável para o armazenamento interno
de instruções para implementações específicas, tais como lógica, sequenciamento, temporização,
contagem e aritmética, para controlar, através de módulos de entradas e saídas, vários tipos de
máquinas ou processos.''}
\end{quote}

Vale destacar quatro elementos dessa definição, porque cada um corresponde a uma decisão de
engenharia:

\begin{enumerate}
  \item \textbf{eletrônico digital} --- não eletromecânico;
  \item \textbf{memória programável} --- a função é definida por programa, não por fiação;
  \item \textbf{funções específicas} --- lógica, sequenciamento, temporização, contagem e
        aritmética: este é o repertório mínimo esperado;
  \item \textbf{módulos de entrada e saída} --- a interface com o processo é parte do produto.
\end{enumerate}

Em uma frase operacional, para o técnico de campo: \emph{CLP é o equipamento que lê o estado das
entradas, executa uma lógica gravada em memória e escreve o resultado nas saídas, repetidamente e em
tempo previsível.}

\index{CLP!definição}

\section{Características gerais}

As características abaixo são o que um comprador espera encontrar em um controlador industrial. O
reconhecimento delas é útil em especificação (Capítulo~\ref{cap:porte}) e em manutenção
(Capítulo~\ref{cap:manutencao}).

\begin{table}[H]
  \centering
  \caption{Características de um controlador programável e sua consequência prática.}
  \label{tab:caracteristicas}
  \begin{tabular}{p{6.2cm} p{8.4cm}}
    \toprule
    Característica & Consequência prática \\
    \midrule
    Operação em ambiente industrial & Suporta faixa de temperatura, vibração, umidade e
      interferência eletromagnética que um computador comum não suporta \\
    Programação e reprogramação rápidas, com mínima interrupção da produção & Altera-se a lógica sem
      refazer painel \\
    Sinalizadores de estado e módulos \emph{plug-in} & Troca de módulo por encaixe, sem mexer na
      fiação dos pontos vizinhos \\
    Dimensão reduzida e baixo consumo & Ocupa menos espaço de painel e dissipa menos calor que o
      painel de relés equivalente \\
    Compatibilidade com diferentes sinais de E/S & Aceita discreto e analógico, em corrente ou
      tensão, com cartões distintos \\
    Capacidade de acionar cargas reais & A saída discreta adequada aciona cargas de alguns ampères,
      como bobinas de contatora e válvulas \\
    Expansão de módulos e de memória & Cresce por módulos, e não por troca de equipamento \\
    Comunicação em rede & Integra-se a outras CPUs, à supervisão e às redes de chão de fábrica \\
    Monitoração do processo por computador & Permite supervisão, controle e aquisição de dados \\
    \bottomrule
  \end{tabular}
\end{table}

\section{Aplicações}

O CLP automatiza processos de sequenciamento, intertravamento, controle de processo e batelada. Além
da manufatura e dos processos contínuos, aparece em automação elétrica, predial e até dentro de
produtos. Os casos mais frequentes são:

\begin{itemize}
  \item máquinas industriais: usinagem, injetoras de plástico, têxteis, calçados;
  \item processos contínuos e de batelada: siderurgia, papel e celulose, petroquímica, química,
        alimentos, mineração;
  \item controle de energia: gestão de demanda e fator de carga;
  \item sinalização, intertravamento e controle PID em utilidades;
  \item aquisição de dados de supervisão em fábricas, edifícios inteligentes e laboratórios;
  \item bancadas de teste automático de componentes;
  \item controle embarcado em máquinas móveis, agrícolas e de construção.
\end{itemize}

\index{aplicações de CLP}

\section{Arquitetura de blocos}

Um CLP é um sistema microprocessado, mais simples que um computador de uso geral, porém projetado
para operar em condições severas. Seus blocos essenciais são a unidade central de processamento
(UCP), as interfaces de entrada e saída e as memórias, todos apoiados por uma fonte de alimentação
tratada e por um barramento interno.

\begin{figure}[H]
  \centering
  \includegraphics[width=\linewidth]{05-arquitetura-clp.pdf}
  \caption{Blocos básicos de um CLP e o caminho do sinal, do campo ao atuador.}
  \label{fig:arquitetura-clp}
\end{figure}

O \textbf{microprocessador} é a inteligência do sistema: obtém os dados dos cartões de entrada, os
processa conforme o programa do usuário e escreve o resultado nos cartões de saída. Um único
processador costuma controlar dezenas ou centenas de pontos de E/S.

A \textbf{fonte de alimentação} recebe tratamento especial de filtragem e estabilização, porque
disponibilidade e qualidade de energia são a diferença entre o ambiente industrial e o escritório. A
alimentação dos \emph{dispositivos de campo}, porém, é externa ao CLP --- veja a caixa de atenção ao
final deste capítulo.

\section{Compacto, modular e com I/O distribuído}

Quanto à estrutura física dos pontos de E/S, há três arranjos:

\begin{table}[H]
  \centering
  \caption{Arranjos físicos de E/S e quando cada um é adequado.}
  \label{tab:arranjos-es}
  \begin{tabular}{p{3.2cm} p{6.0cm} p{5.4cm}}
    \toprule
    Arranjo & Como é & Quando convém \\
    \midrule
    Compacto
      & Todos os componentes em um único módulo, exceto o terminal de programação; costuma usar
        microcontrolador no lugar do conjunto processador, memória e interface
      & Aplicações pequenas, poucos pontos, custo e espaço reduzidos \\
    Modular
      & Um ou mais racks, com fonte, CPU e cartões de E/S encaixáveis, escolhidos conforme a
        necessidade
      & Aplicações que crescem, exigem cartões especiais ou troca a quente \\
    I/O distribuído
      & Ao menos uma CPU, fonte e módulos de E/S separados, comunicando por rede de campo
        (\emph{fieldbus})
      & Instrumentação espalhada por área extensa, com economia de cabeamento \\
    \bottomrule
  \end{tabular}
\end{table}

Na prática os arranjos se misturam: existem controladores com E/S embutida que também aceitam cartões
modulares e ainda operam com E/S distribuída por \emph{fieldbus}.

\section{O que distingue o CLP de um computador comum}

A estrutura é parecida, o propósito não é:

\begin{itemize}
  \item \textbf{determinismo} --- o ciclo se repete em tempo previsível; no computador comum, o
        escalonamento do sistema operacional não garante prazo;
  \item \textbf{fonte e entradas tratadas} --- filtragem, estabilização e imunidade a ruído são
        requisito de projeto, não acessório;
  \item \textbf{invólucro e faixa de operação} --- temperatura, vibração e umidade de ambiente
        industrial;
  \item \textbf{diagnóstico} --- sinalização de estado por ponto e por módulo, pensada para quem
        mantém a máquina.
\end{itemize}

\begin{atencao}
A alimentação dos dispositivos de campo deve ser fornecida externamente. A fonte do CLP é
dimensionada para a parte interna do equipamento, não para alimentar sensores, válvulas e
sinalizadores. Tentar alimentar carga de campo pela fonte interna é uma das causas mais comuns de
queima de fonte e de reset indevido.
\end{atencao}

\begin{pratica}
Abra o catálogo de um controlador disponível na sua escola e localize, para ele:
(a) o arranjo de E/S (compacto, modular ou distribuído);
(b) a tensão de alimentação da CPU e a das saídas digitais;
(c) o tempo de varredura por mil instruções;
(d) quais cartões de entrada analógica existem e com quantos bits de resolução.
Relacione cada dado com uma linha da Tabela~\ref{tab:caracteristicas}.
\end{pratica}

\begin{exercicios}
\begin{enumerate}[label=\alph*)]
  \item Reescreva a definição da NEMA em linguagem de manutenção, em no máximo três linhas.
  \item Um equipamento lê entradas, executa lógica e aciona saídas, mas não resiste a vibração e
        calor de painel. Ele é um CLP? Justifique pela definição.
  \item Explique por que o CLP compacto costuma empregar microcontrolador em vez do conjunto
        processador, memória e interface.
  \item Em que situação o I/O distribuído é mais barato que o modular, e em que situação é mais caro?
  \item Um aluno afirma que ``um CLP é um computador industrial''. Aponte duas diferenças que
        tornam essa afirmação tecnicamente imprecisa.
\end{enumerate}
\end{exercicios}
